\chapter{Applications, évaluation et discussion}
\label{ch:evaluation}

% Pour un travail théorique, utilisez ce chapitre pour les exemples,
% conséquences ou limites. Pour un travail empirique, séparez protocole,
% résultats et discussion.

Considérons le graphe de la \cref{fig:example-graph}. Ses arêtes directes sont
$a\to b$, $b\to c$, $c\to d$ et $a\to d$. La figure est dessinée avec TikZ :
elle reste nette quel que soit le zoom.

\begin{figure}[htbp]
  \centering
  \begin{tikzpicture}[
  vertex/.style={circle,draw=LMFINavy,thick,minimum size=8mm,inner sep=0pt},
  edge/.style={-{Stealth[length=2.4mm]},thick,LMFITeal},
  node distance=2.2cm
]
  \node[vertex] (a) {$a$};
  \node[vertex, right=of a] (b) {$b$};
  \node[vertex, right=of b] (c) {$c$};
  \node[vertex, below=1.0cm of c] (d) {$d$};

  \draw[edge] (a) -- (b);
  \draw[edge] (b) -- (c);
  \draw[edge] (c) -- (d);
  \draw[edge] (a) to[bend right=18] (d);
\end{tikzpicture}

  \caption{Graphe orienté à quatre sommets utilisé pour évaluer l'accessibilité.}
  \label{fig:example-graph}
\end{figure}

L'application de l'\cref{alg:warshall} donne la matrice réflexive d'accessibilité
du \cref{tab:reachability}. Le sommet $a$ atteint tous les sommets, tandis que
$d$ n'atteint que lui-même. L'entrée de la ligne $b$ et de la colonne $d$
représente le chemin $b\to c\to d$, bien qu'il n'existe pas d'arête directe
$b\to d$.

\begin{table}[htbp]
  \centering
  \caption{Matrice réflexive d'accessibilité de la \cref{fig:example-graph}.}
  \label{tab:reachability}
  \begin{tabular}{c *{4}{>{\centering\arraybackslash}p{1.2cm}}}
    \toprule
    De $\backslash$ vers & $a$ & $b$ & $c$ & $d$ \\
    \midrule
    $a$ & 1 & 1 & 1 & 1 \\
    $b$ & 0 & 1 & 1 & 1 \\
    $c$ & 0 & 0 & 1 & 1 \\
    $d$ & 0 & 0 & 0 & 1 \\
    \bottomrule
  \end{tabular}
\end{table}

Cet exemple vérifie la cohérence interne plutôt que la performance. Un véritable
chapitre expérimental doit préciser les données, la version du logiciel,
la métrique et l'incertitude, et discuter les cas négatifs : l'accessibilité
seule perd les longueurs de chemins.
